\section{Monotonicity and the first derivative}

The derivative measures local change. This makes it a powerful tool for understanding whether a function is increasing or decreasing. The sign of the derivative tells us the local direction of change.

This is the first major use of derivatives in function analysis.

\subsection*{Positive and negative derivative}

If \(f'(x)>0\) on an interval, then the function is increasing on that interval. If \(f'(x)<0\), then the function is decreasing.

\begin{theorembox}
Let \(f\) be differentiable on an interval.
\begin{itemize}
\item If \(f'(x)>0\) for all \(x\) in the interval, then \(f\) is increasing there.
\item If \(f'(x)<0\) for all \(x\) in the interval, then \(f\) is decreasing there.
\end{itemize}
\end{theorembox}

This theorem should be read geometrically. A positive derivative means that tangent lines have positive slope. The graph rises as we move from left to right. A negative derivative means that tangent lines have negative slope. The graph falls.

\subsection*{A basic example}

\begin{examplebox}
Study the monotonicity of
\[
f(x)=x^2.
\]
The derivative is
\[
f'(x)=2x.
\]
The sign of \(2x\) depends on \(x\):
\[
2x<0 \quad\text{for } x<0,
\]
\[
2x=0 \quad\text{for } x=0,
\]
\[
2x>0 \quad\text{for } x>0.
\]
Therefore \(f\) is decreasing on \((-\infty,0)\) and increasing on \((0,\infty)\).
\end{examplebox}

This conclusion matches the graph of the parabola. The derivative gives an algebraic way to detect what the graph shows visually.

\subsection*{Sign charts}

A sign chart is a simple way to organize where the derivative is positive, negative, or zero. The important points are usually where the derivative is zero or undefined. These points divide the domain into intervals. On each interval, we check the sign of the derivative.

\begin{examplebox}
Study the monotonicity of
\[
f(x)=x^3-3x.
\]
First compute the derivative:
\[
f'(x)=3x^2-3=3(x^2-1)=3(x-1)(x+1).
\]
The derivative is zero at
\[
x=-1
\qquad\text{and}\qquad
x=1.
\]
These points divide the real line into three intervals:
\[
(-\infty,-1),
\qquad
(-1,1),
\qquad
(1,\infty).
\]
Check the sign of \(f'(x)=3(x-1)(x+1)\):
\begin{itemize}
\item on \((-\infty,-1)\), both factors are negative, so \(f'(x)>0\),
\item on \((-1,1)\), one factor is negative and one is positive, so \(f'(x)<0\),
\item on \((1,\infty)\), both factors are positive, so \(f'(x)>0\).
\end{itemize}
Therefore \(f\) is increasing on \((-\infty,-1)\), decreasing on \((-1,1)\), and increasing on \((1,\infty)\).
\end{examplebox}

The derivative does not only compute slopes at isolated points. It organizes the whole graph into intervals of increasing and decreasing behavior.

\subsection*{Where the derivative is zero}

A derivative equal to zero means that the tangent line is horizontal. This may indicate a local maximum or local minimum, but it does not always do so.

\begin{examplebox}
For
\[
f(x)=x^3,
\]
we have
\[
f'(x)=3x^2.
\]
Thus
\[
f'(0)=0.
\]
However, \(x^3\) is increasing on both sides of \(0\). The point \(x=0\) is not a local maximum or a local minimum.
\end{examplebox}

This example is important. A zero derivative is a signal to investigate, not an automatic conclusion.

\begin{summarybox}
When using the first derivative for monotonicity, ask:
\begin{itemize}
\item What is the domain of the function?
\item Where is \(f'(x)=0\)?
\item Where is \(f'(x)\) undefined?
\item What intervals are created by these points?
\item What is the sign of \(f'(x)\) on each interval?
\item What does that sign say about the graph?
\end{itemize}
\end{summarybox}

The first derivative is a language for reading growth and decline.